\section{偏导数为何不等于可微}
\label{sec:02-Calculus/06-Multivariable-Calculus/02}

\subsection{两个坐标方向的导数都算出来了，然后呢}

你有一个二元函数 \(f(x,y)\)。在某个点 \((a,b)\)，你把 \(y\) 固定为 \(b\)，只动 \(x\)，算出一个导数 \(f_x(a,b)\)；再把 \(x\) 固定为 \(a\)，只动 \(y\)，算出 \(f_y(a,b)\)。两个偏导数都存在，数值也很好看。

你的第一反应可能是：这函数在该点``可导''，局部线性近似的切平面就是
\[
z\approx f(a,b)+f_x(a,b)(x-a)+f_y(a,b)(y-b).
\]
毕竟单变量微积分里，导数存在就已经几乎等于可微了。拿一个看起来温和的函数试试。

定义
\[
f(x,y)=
\begin{cases}
\dfrac{xy}{\sqrt{x^2+y^2}},&(x,y)\ne(0,0),\\
0,&(x,y)=(0,0).
\end{cases}
\]

沿 \(x\) 轴：函数恒为零，所以 \(f_x(0,0)=0\)。沿 \(y\) 轴：同样恒为零，\(f_y(0,0)=0\)。两个偏导都漂漂亮亮地存在。如果偏导存在就足以保证可微，那么候选的切平面就是 \(z=0\)，而近似误差应该满足
\[
\frac{|f(x,y)-0|}{\sqrt{x^2+y^2}}\to0.
\]

沿 \(y=x\) 走一步看看：
\[
\frac{|f(x,x)|}{\sqrt{x^2+x^2}}
=\frac{|x^2/\sqrt{2x^2}|}{\sqrt{2}\,|x|}
=\frac{|x|/\sqrt{2}}{\sqrt{2}\,|x|}
=\frac12.
\]
误差没有趋零。不论步长取多小，相对误差稳稳停在 \(1/2\)。偏导数都在，切平面却不存在。

\subsection{偏导数只沿坐标轴切片，可微性需要所有方向}

偏导数的定义本身就是坐标轴上的极限：
\[
f_x(a,b)=\lim_{h\to0}\frac{f(a+h,b)-f(a,b)}h,
\qquad
f_y(a,b)=\lim_{h\to0}\frac{f(a,b+h)-f(a,b)}h.
\]

固定其他变量、只动一个变量的操作，几何上是在曲面 \(z=f(x,y)\) 上切两刀——一刀平行于 \(xz\) 平面，一刀平行于 \(yz\) 平面。你得到了两条交线在切点处的切线斜率。

计算偏导数时把其余变量视为常数，这本身就是单变量求导的套路。例如
\[
f(x,y)=x^2y+e^{xy}
\]
有
\[
f_x=2xy+ye^{xy},\qquad
f_y=x^2+xe^{xy}.
\]
按公式走完全没问题。

但坐标系里只有两个坐标轴方向，而平面上的方向有无穷多个。两个方向的斜率意味着你有了一个候选的切平面——过该点、沿 \(x\) 方向斜率 \(f_x\)、沿 \(y\) 方向斜率 \(f_y\) 的那个平面。这个平面在其他方向上的斜率由线性结构自动决定。问题是，函数在那些非坐标方向上的实际变化，跟平面预测的变化一致吗？

\subsection{一锅汤：两个坐标方向不能代替一条斜线}

调汤时每次只改一种调料，可以估计盐量和水量分别对味道的影响。设 \(x\) 是盐量偏离配方的值，\(y\) 是水量偏离配方的值，\(f(x,y)\) 是味道评分相对于配方点的偏差。按偏导数的方法，固定水量、只尝不同盐量，得到一个关于盐的边际变化 \(f_x\)；固定盐量、只尝不同水量，得到关于水的边际变化 \(f_y\)。

这两个边际信息合起来，只覆盖了坐标轴上从配方点出发的两条线。别人的配方可能同时增减盐和水：\(x\) 和 \(y\) 以某种比例联合变化。沿 \(y=x\) 的方向，涉及盐和水的相互补偿——浓度由比值决定。这个方向上的实际味道变化，是仅凭两条坐标轴上的切片信息无法推断的。

前面 \(xy/\sqrt{x^2+y^2}\) 的例子把这个盲点做绝了：坐标轴上没有任何一阶变化——\(f_x=f_y=0\)，函数沿坐标轴保持恒零。但沿 \(y=x\) 联合移动时，函数值与步长同阶增长，而切平面预测的却是零增长。一次只改一个因素可以估算偏导，却不能单独证明一个线性模型对所有联合改动都有效。

\subsection{可微性的准确定义：存在一个线性映射}

\(f:\R^n\to\R^m\) 在 \(\mathbf a\) 可微，若存在一个线性映射 \(A\) 使
\[
f(\mathbf a+\mathbf h)
=f(\mathbf a)+A\mathbf h+\mathbf r(\mathbf h),
\]
且
\[
\frac{\|\mathbf r(\mathbf h)\|}{\|\mathbf h\|}\to0.
\]

\(A=Df(\mathbf a)\) 称为全导数。它不是某个方向上的斜率，而是一个线性算子——输入任意方向的小位移 \(\mathbf h\)，输出函数值的一阶变化量。误差项 \(\mathbf r(\mathbf h)\) 相对于步长 \(\|\mathbf h\|\) 更高阶地小：你把步长缩小一半，误差缩小到原来的四分之一以下。

这条定义没有提到任何偏导数。它直接要求存在一个线性近似，对所有方向同时有效。偏导数是可微性的推论，不是前提。

当输出是标量时，
\[
Df(\mathbf a)\mathbf h
=\nabla f(\mathbf a)^{\trans}\mathbf h,
\]
其中 \(\nabla f\) 是把各偏导排列成向量的梯度记号。当 \(f:\R^n\to\R^m\) 时，线性映射由 Jacobian 矩阵表示。梯度与 Jacobian 在下一节正式定义并展开。

\subsection{切平面来自全微分}

对二元标量函数，在 \((a,b)\) 可微时，
\[
f(a+\Delta x,b+\Delta y)
\approx f(a,b)+f_x(a,b)\Delta x+f_y(a,b)\Delta y.
\]
图像 \(z=f(x,y)\) 在 \((a,b,f(a,b))\) 处的切平面为
\[
z=f(a,b)+f_x(a,b)(x-a)+f_y(a,b)(y-b).
\]

例如 \(f(x,y)=x^2+xy\) 在 \((1,2)\)：\(f(1,2)=3\)，\(f_x=2x+y\) 给出 \(f_x(1,2)=4\)，\(f_y=x\) 给出 \(f_y(1,2)=1\)，切平面为
\[
z=3+4(x-1)+(y-2).
\]

这张平面做的事情是预测小输入变化引起的输出变化——在 \((1,2)\) 附近，\(x\) 每增加一个单位，输出约增 4；\(y\) 每增加一个单位，输出约增 1。它不是在宣称曲面全局是平面，而是一个局部的一阶模型。离切点越远，预测就越不可靠。

\subsection{一个实用的检查路径}

若各一阶偏导数在点附近存在并连续（即函数是 \(C^1\) 的），则函数在该点可微。这是最常用的充分条件——检查偏导数的连续性远比直接验证可微定义中的误差项来得方便。

但它只是充分条件，不是必要条件。存在这样的函数：偏导数都在，可微性也成立，但偏导数本身不连续。例如
\[
f(x,y)=
\begin{cases}
(x^2+y^2)\sin\frac{1}{\sqrt{x^2+y^2}},&(x,y)\ne(0,0),\\
0,&(x,y)=(0,0)
\end{cases}
\]
在原点可微（因为函数值与距离平方同阶），但偏导数在原点附近剧烈振荡。

逻辑箭头应该印在脑海里：
\[
\text{一阶偏导附近连续}\;(C^1)
\;\Longrightarrow\;
\text{可微}
\;\Longrightarrow\;
\text{连续且各偏导存在}.
\]

两个箭头都不可逆。上一节的反例 \(xy/\sqrt{x^2+y^2}\) 已经证明``偏导存在 + 连续 \(\nRightarrow\) 可微''；本节末尾的例子证明了``可微 \(\nRightarrow\) \(C^1\)''。

看待这三个层次的一个实际习惯：先算偏导并检查它们是否在原点的某个邻域内连续——如果是，可微直接到手，不需要多费功夫。如果偏导不连续，不要立刻断定不可微；回到定义，考察误差项 \(\mathbf r(\mathbf h)/\|\mathbf h\|\) 是否真趋于零。大部分``看起来病态''但实际可微的函数，都是在原点附近以振荡的方次抵消了误差。

下一节把可微性与方向导数对接起来：有了可微性做担保，梯度向量就成为了所有方向的``变化率编码器''。
